\section*{Chapter \ref{ch:m3:agriculturals-and-softs} --- Agriculturals and Softs}

\begin{solution}{exo:m3:agriculturals-and-softs:1}
$1.2/14.8 = 8.1\%$.
\end{solution}

\begin{solution}{exo:m3:agriculturals-and-softs:2}
At \$4.00: 7\% is 28 cents, rounded to 30; expanded 45. At \$6.00: 42 cents, rounded to 40; expanded 60.
\end{solution}

\begin{solution}{exo:m3:agriculturals-and-softs:3}
July delivers before the US harvest, so only grain already in store can be delivered: \omterm{def:m3:agriculturals-and-softs:cropyear}{old crop}. December
delivers after the harvest: \omterm{def:m3:agriculturals-and-softs:cropyear}{new crop}.
\end{solution}

\begin{solution}{exo:m3:agriculturals-and-softs:4}
$300 \times 0.022 + 50 \times 0.11 - 10.00 = 6.60 + 5.50 - 10.00 = \$2.10$ a bushel.
\end{solution}

\begin{solution}{exo:m3:agriculturals-and-softs:5}
One locked day (30 cents), then the remaining 30 cents is within the \omterm{def:m3:agriculturals-and-softs:limit}{expanded limit} and the market trades
at the fair price on day 2. The loss is 60 cents, \$3{,}000 per contract.
\end{solution}

\begin{solution}{exo:m3:agriculturals-and-softs:6}
Crowded longs must sell on the same news, and there are few new buyers left: forced selling adds to the
fundamental move, and in a limit market it can lock the price for days.
\end{solution}

\begin{solution}{exo:m3:agriculturals-and-softs:7}
Net long 24.4\% of open interest on 22 March 2022; net short 22.8\% on 9 July 2024; net short in 300 of 559
weeks.
\end{solution}

\begin{solution}{exo:m3:agriculturals-and-softs:8}
Correlation does not give direction: funds that follow trends buy after prices rise, so their positions
may follow prices rather than lead them. The report is also published three days after the positions it
shows. A strategy needs evidence that positions predict later returns, tested out of sample on data
available at the time.
\end{solution}

\begin{solution}{pb:m3:agriculturals-and-softs:1}
\textbf{1.} 30 and 45 cents. \textbf{2.} \$50 per cent per contract (\$12.50 per quarter-cent tick).
\textbf{3.} \$4.20, locked. \textbf{4.} \$3.75, locked at the \omterm{def:m3:agriculturals-and-softs:limit}{expanded limit}. \textbf{5.} Day 3, at \$3.70.
\textbf{6.} $50 \times 0.30 \times 5\,000 = \$75{,}000$. \textbf{7.} \$112{,}500. \textbf{8.} \$187{,}500
(\$3{,}750 per contract). \textbf{9.} \$200{,}000. \textbf{10.} Nothing during the lock: it becomes a
market order that fills only when trading resumes, near \$3.70. \textbf{11.} Options may keep trading and
price the fair value, at a premium reflecting it; synthetic futures from options reveal the locked
contract's true price. \textbf{12.} Possibly other months or an overseas contract, with \omterm{def:m3:physical-commodity-markets:basisrisk}{basis risk}.
\textbf{13.} Its short hedge gains margin in cash while its crop loses value: it is paid, not called.
\textbf{14.} Limits give time to raise margin and slow panics; the cost is illiquidity. \textbf{15.} So that
the margin for two or three locked days is covered and the loss is tolerable, not by the usual volatility.
\textbf{16.} How crowded the positioning is: crowded positions add forced selling. \textbf{17.} It caps the
daily move, not the total one: the price keeps moving on later days. \textbf{18.} Liquidity is worth most
when you cannot trade: the ability to hedge elsewhere or to fund margin is what survives the lock.
\textbf{19.} \emph{Named result:} \$3{,}750 per contract, \$187{,}500 for 50 contracts. \textbf{20.} A
maximum daily move that delays, but does not cancel, a price change.
\end{solution}

\begin{solution}{iq:m3:agriculturals-and-softs:1}
Ending stocks over total use: the buffer carried into the next harvest. When the buffer is thin, only a
higher price can ration use, and demand for a staple is inelastic, so the price rises steeply as the ratio
falls; when it is ample, price is set by the cost of storage.
\iqlookfor{inelastic demand and a convex price response.}
\end{solution}

\begin{solution}{iq:m3:agriculturals-and-softs:2}
A market at its daily limit with orders on one side only: positions cannot be closed. Manage it by sizing
for several locked days, by hedging in instruments that keep trading (options, other contracts), and by
holding margin liquidity.
\iqlookfor{size to the lock, alternatives, funding.}
\end{solution}

\begin{solution}{iq:m3:agriculturals-and-softs:3}
Use the release date (Friday), not the report date (Tuesday), as the time the data are known; normalise
positions (share of open interest, z-score against the past only); test on out-of-sample years with
non-overlapping returns; account for revisions and category changes.
\iqlookfor{release-date alignment and honest normalisation.}
\end{solution}

\begin{solution}{iq:m3:agriculturals-and-softs:4}
Beans in, meal and oil out; its margin is the crush. It hedges by buying soybean futures and selling meal
and oil futures in the board ratio, locking the margin on its planned throughput.
\iqlookfor{the product-input spread and the ratio.}
\end{solution}

\begin{solution}{iq:m3:agriculturals-and-softs:5}
Daily returns are truncated at the limit, so a one-day VaR understates risk. Use multi-day scenarios with
the remaining move carried forward, or measure over the horizon needed to exit; include \omterm{def:m3:agriculturals-and-softs:limit}{expanded limits}.
\iqlookfor{the liquidity horizon, not the one-day truncated distribution.}
\end{solution}

\begin{solution}{iq:m3:agriculturals-and-softs:6}
Inputs: consensus and prior estimates, positions and limits, liquidity by contract. Before noon: set risk
limits, cancel resting orders that should not fill on news, pre-compute scenarios. At noon: parse the
release, compute surprises, route orders with price protection, switch to options or other months if a
contract locks; afterwards, reconcile and report.
\iqlookfor{parsing, surprise computation, protection and fallbacks.}
\end{solution}
